%%
\documentclass[manuscript,screen,12pt, nonacm]{acmart}
\let\Bbbk\relax % Fix for amssymb clash 
\usepackage{standalone}
\usepackage{vmlmacros}
\AtBeginDocument{%
  \providecommand\BibTeX{{%
    \normalfont B\kern-0.5em{\scshape i\kern-0.25em b}\kern-0.8em\TeX}}}
\usepackage{outlines}
%~\usepackage{natbib}
\setlength{\headheight}{14.0pt}
\setlength{\footskip}{13.3pt}
\citestyle{acmauthoryear}
\begin{document}
\title{An Alternative to Pattern Matching, Inspired by Verse}
\author{Roger Burtonpatel}
\email{roger.burtonpatel@tufts.edu}
\affiliation{%
  \institution{Tufts University}
  \streetaddress{419 Boston Ave}
  \city{Medford}
  \state{Massachusetts}
  \country{USA}
  \postcode{02155}
}
\bibliographystyle{plainnat}
% \renewcommand\refname{}
% ~\setcitestyle{sort&compress}


%~\author{Norman Ramsey}
%~\email{nr@cs.tufts.edu}
%~\affiliation{%
%~\institution{Tufts University}
%~\streetaddress{177 College Ave}
%~\city{Medford}
%~\state{Massachusetts}
%~\country{USA}
%~\postcode{02155}
% }

%~\author{Milod Kazerounian}
%~\email{milod.mazerouniantufts.edu}
%~\affiliation{%
%~\institution{Tufts University}
%~\streetaddress{177 College Ave}
%~\city{Medford}
%~\state{Massachusetts}
%~\country{USA}
%~\postcode{02155}
% }

\renewcommand{\shortauthors}{Burtonpatel}

% Pretty dece: have Steve read. 
\begin{abstract}
    Pattern matching % S
    appeals % V
    to functional programmers for its expressiveness and good cost model, 
    but 
    it % S
    fails to express % V
    certain computations without verbosity. Verbosity % S 
    can be reduced % V
    by using extended forms of pattern matching,
    but extensions % (S = the problem)
    are % V
    not standardized, and making certain extensions like or-patterns and pattern
    guards cooperate in a large-scale programming language is a significant
    implementation challenge that has not seen success until recent years. 
    
    % the most desirable combination of extensions could not
    % until recently found in any one popular programming language.
    
    Alternatively to pattern matching, computations can be expressed succinctly
    by using equations, as has been demonstrated by the Verse programming
    language. But such computations may have time and space costs that can be
    hard to predict. 
    % But Verse's cost model % S-sub
    % is % V-sub
    % a challenge because of backtracking and multiple
    % results. 
    As a compromise, 
    we % S
    ~propose % V
    a new language,~\VMinus, % object 
    which~uses equations in a limited way that makes time and space costs easy
    to predict. 
    % to succinctly express computation while
    % enjoying a friendlier cost model by prohibiting backtracking and multiple
    % results. 
    In
    comparative examples,
    \VMinus~ % S 
    expresses % V
    computations as succinctly as pattern matching with popular extensions, and just like pattern matching, it % S
    can be compiled % V 
    to a decision tree. 
    
    % Evaluation
    To evaluate % v
    the succinctness of \VMinus,
    we compare % V
    it with a novel pattern-matching language, \PPlus. \PPlus
    demonstrates \rab{want better verb} % V
    three popular pattern matching extensions 
    in unison: side conditions, or-patterns, and pattern guards. 
    We evaluate % V 
    the expressiveness of both languages using four smaller contrived examples 
    and a larger example from \_. 

    An implementation of~\VMinus, \PPlus, and the match compiler can be found
    at~\url{https://github.com/rogerburtonpatel/vml}.
    % Question: implied that I've implemented D here, or not? 
    %%%%%%% Old, with "I show":
    % I % S
    % ~show % V
    %  that~\VMinus's equations capture the desirable properties of
    % pattern matching with popular extensions, 
    % and I % S-sub
    % ~also
    %  show % V 
    % that~\VMinus~can be compiled to a decision tree and therefore has
    % the same desirable cost model as pattern matching. Finally, 
    % I % S 
    % ~have % V 
    % implemented~\VMinus in Standard ML, alongside an algorithm for compiling it
    % to a decision tree. 
\end{abstract}

\maketitle

\section{Introduction}
% Subjects: Pattern matching and Equations 

Perhaps the most beloved tool among functional programmers for examining and
deconstructing data is pattern matching. 
% It does so implicitly by matching constructed data
% directly against a number of possible forms.
% Maybe combine? 
Pattern matching is also an established and well-researched topic
\citep{wadler1987views, macqueen1985tree, burton1993pattern, palao1996new,
maranget, bpc}. It is appreciated by programmers and researchers alike for two
main reasons: It enables~\it{implicit} data deconstruction, and it has a
desirable cost model. Specifically (regarding the latter), pattern matching can
be compiled to a~\it{decision tree}, a data structure that enforces linear
runtime performance by guaranteeing no part of the data will be examined more
than once.~\citep{maranget}

However, pattern matching cannot express certain common computations succinctly,
forcing programmers who wish to express these computations to duplicate code,
nest~\it{case} expressions, and create multiple points of truth. To mitigate
this, designers of popular programming languages have introduced~\it{extensions}
to pattern matching (Section~\ref{extensions}). 

Extensions strengthen pattern matching, but they are not standardized, so each
popular programming language with pattern matching features its own unique suite
of extensions. Extensions are subject to the discretion of the individual
language designer, not ubiquitous. Rather than continuing to extend pattern
matching~\it{ad hoc}, a worthwhile goal could be to find an alternative that
doesn't need extensions. A tempting possibility was introduced last year by the
programming language Verse~\citep{verse}. In Verse, a programmer can deconstruct
data using a different tool the language offers:~\it{equations}. Equations are
expressive and flexible, and it appears that they can express everything that
pattern matching can, including with popular extensions. 

But a full implementation of Verse is complicated, cost-wise. Verse is a
functional logic programming language, and expressions can backtrack at runtime
and return multiple results, both of which are hard to predict in their costs. 

% A worthwhile goal could be to harmonize the expressive quality of Verse's
% equations with the decision tree property of patterns.

\bf{In this thesis, I~show} that the expressive quality of Verse's equations and
the decision-tree property of patterns can be combined in a single language.
Since the language is a streamlined adaptation of Verse with a reduced feature
set, I~call it~\VMinus (“V minus”). 
% I~also show that~\VMinus subsumes pattern matching with popular extensions. 

To support this claim, I~have formalized~\VMinus with a big-step operational
semantics (Section~\ref{vminus}), I~have formalized decision trees into a core
language~\D (“D”) with a big-step operational semantics (Section~\ref{d}), 
% I~have formalized pattern matching in a core language~\PPlus (“P plus”) with a
% big-step operational semantics (Section~\ref{pplus}),
I~have formalized a translation from~\VMinus to~\D (Sections~\ref{vminustod}),
and I~have implemented both languages in Standard ML. 



\section{Pattern Matching and its Extensions}
\label{pmandextensions}
    \textbf{\rab{GOAL: Introduce pattern matching in a way that's neither
    dismissive of the interesting problem nor insulting to the audience's
    intelligence and knowledge of functional programming.}}


    Pattern matching is a battle-tested and beloved features in declarative and
    functional programming, but it often requires extensions to simplify cases
    that are otherwise troublesome. 
    
    
    Specifically, in sophisticated programs that involve complex
    decision-making, pattern matching without extensions tends to duplicate
    code. Consider this example of "simple" pattern matching (without
    extensions): 

    \rab{Example; will discuss which to use. Adaptation of 'cellar' example
    seems good}.

    \rab{Example says along the lines of: "Without extensions to pattern
    matching, expressing the algorithm looks like this: \dots}. 

    The code gets the job done, but it is cumbersome to write and read because
    of the large amount of duplication, which arises from the limitations of
    what can be expressed in simple pattern matching. 
    
    Furthermore, if the spec changes, will the programmer be shrewd enough to
    find and modify all the duplicated parts of the program to match the change?
    Such an example demands an improvement. Indeed, we can do better, by 
    augmenting pattern matching with \textit{extensions}. 


    \subsection{Popular extensions to pattern matching}
    \label{extensions}

        % Extensions to pattern matching simplify cases that are otherwise
    % troublesome. In this section, we~illustrate several such cases. The three
    % extensions we~describe are those commonly found in the literature and
    % implemented in compilers: side conditions, pattern guards, and or-patterns. 
    
    % To denote pattern matching~\it{without} extensions, we~coin the
    % term~\it{simple pattern matching}. In bare pattern matching, a pattern has
    % one of two syntactic forms: a name or an application (of a value constructor
    % to zero or more patterns). 


    \rab{TODO: Format this paragraph so it presents information as clearly as possible.}
    Duplicated code, especially in a complex case expression, can make code
    difficult to maintain and more bug-prone. To address this, many large-scale
    functional programming languages have added extensions to pattern matching,
    most commonly \textit{side conditions}(cite), which extend a pattern match with a
    boolean guard, \textit{or-patterns}(cite), which allow multiple branches in a case
    expression to share a right-hand side, and \textit{pattern guards}(cite), which
    allow multiple sequential pattern-matches and computations to guard a
    right-hand side. 

    \rab{Side condition example}
    Side conditions do \dots which eliminate the need to \dots, 
    illustrated in this small example of measuring shapes: 

        \rab{The examples are funny, but they should be more serious for this 
        submission.}

        \rab{TODO: merge figures, make them appear in the right place}
        \begin{figure}[ht]
        \begin{verbatim}
            type standard_shape = Square of float 
                    | Triangle of float * float
                    | Trapezoid of float * float * float
                    

            let exclaimTall sh =
            match sh with 
              | Square s -> if s > 100.0 
                            then "Wow! That's a sizeable square!"
                            else "Your shape is small." 
              | Triangle (_, h) -> 
                            if h > 100.0 
                            then "Goodness! Towering triangle!"
                            else "Your shape is small." 
              | Trapezoid (_, _, h) -> 
                            if h > 100.0
                            then "Zounds! Tremendous trapezoid!"
                            else "Your shape is small." 
            \end{verbatim}    
        \caption{An invented function~\tt{exclaimTall} in OCaml combines pattern
        matching with an~\tt{if} expression, and is not very pretty.}   
        \Description{An implementation of a function exclaimTall in OCaml, which
        uses pattern matching and an if statement.}
        \label{fig:ifexclaimtall}
    \end{figure}


       \begin{figure}[]
            \begin{verbatim}
                let exclaimTall sh =
                  match sh with 
                    | Square s when s > 100.0 ->
                            "Wow! That's a sizeable square!"
                    | Triangle (_, h) when h > 100.0 ->
                            "Goodness! Towering triangle!"
                    | Trapezoid (_, _, h) when h > 100.0 -> 
                            "Zounds! Tremendous trapezoid!"
                    | _ ->  "Your shape is small." 
                \end{verbatim}
            \caption{With a side condition,~\tt{exclaimTall} in OCaml becomes
            simpler and more adherent to the Nice~Properties.} 
            \Description{An implementation of a function exclaimTall in OCaml,
            which uses pattern matching and a side condition.}
            \label{fig:whenexclaimtall}
        \end{figure}



    \rab{Or pattern example}
    Or patterns do \dots which eliminate the need to \dots

     \begin{figure}
            \begin{center}
                \begin{verbatim}
        let game_of_token token = match token with 
            | BattlePass f      -> ("Fortnite", f)
            | ChugJug f         -> ("Fortnite", f)
            | TomatoTown f      -> ("Fortnite", f)
            | HuntersMark f     -> ("Bloodborne", 2 * f)
            | SawCleaver f      -> ("Bloodborne", 2 * f)
            | MoghLordOfBlood f -> ("Elden Ring", 3 * f)
            | PreatorRykard f   -> ("Elden Ring", 3 * f)
            | _                 -> ("Irrelevant", 0)
                \end{verbatim}
            \end{center}    

        \caption{\tt{game\_of\_token}, with redundant right-hand sides,
        should raise a red flag.} 
        \Description{An implementation of game_of_token using bare patterns.}
        \label{fig:baregot}
        \end{figure}


          \begin{figure}
    \begin{center}
    \begin{verbatim}
let game_of_token token = match token with 
    | BattlePass f | ChugJug f | TomatoTown f  -> ("Fortnite", f)
    | HuntersMark f | SawCleaver f             -> ("Bloodborne", 2 * f)
    | MoghLordOfBlood f | PreatorRykard f      -> ("Elden Ring", 3 * f)
    | _                                        -> ("Irrelevant", 0)
    \end{verbatim}
    \end{center}    
    \caption{Or-patterns condense~\tt{game\_of\_token}
    significantly, and it is easier to read line-by-line.} 
    \Description{An
    implementation of game_of_token using or-patterns.}
    \label{fig:orgot}
    \end{figure}

    \newpage
    Individually, each of these three extensions grants a powerful abstraction
    that reduces code duplication. It follows that they should retain their
    powers of brevity when operating in unison, and indeed this is the case:
    with these three extensions co-operating, the previous example reduces like
    this \dots

    Given the expressive power of this combination and each extension's
    individual popularity, it would seem obvious that they are all supported by
    myriad functional programming languages. But that's not true: among Haskell,
    OCaml, Scala, Erlang/Elixir, Rust, F\#, Rocq, and Agda~\citep{haskell,
    ocaml, scala, erlang, elixir, rust, fsharp, agda}, only Haskell has all
    three of side conditions, or patterns, and pattern guards- and only as of
    2023! In fact, when the original idea for this work was conceived in 2023, a
    primary motivator for suggesting an alternative to pattern matching was that
    there was \textit{no} language with all three of them. 

    - No strict language with all three + significant implementation effort in
    Haskell to get them all working -> evidence enough that having all of them
    is a non-trivial problem. 
    
    We~find the extension story somewhat unsatisfying. At the very least,
    we~want to be able to use pattern matching, with the extensions we~want, in
    a single language. Or, we~want an alternative that gives me the expressive
    power of pattern matching with these extensions. 

        \section{Equations}

    An intriguing alternative to pattern matching exists in~\it{equations}, from
    the Verse Calculus (\VC), a core calculus for the functional logic
    programming language~\it{Verse}~\citep{antoy2010functional,
    hanus2013functional, verse}. (For the remainder of this work, we~use “\VC”
    and “Verse” interchangeably.)

    Verse's equations and choice mechanism allow for much of the same expressive
    power as pattern matching with extensions. Rather than go into too much
    detail (the curious may examine the Verse paper, which is well worth a
    read), we support this claim with a familiar example: 



    The brevity of the code in the figures is is promising for~\VC. If it rivals
    pattern matching with popular extensions in expressiveness, and~\VC
    does everything using only equations and choice, it seems like the language
    is an intriguing option for writing code! 


    But the VC semantics has a cost model that is hard to predict. Why? In~\VC,
    names (logical variables) are~\it{values}, and they can just as easily be
    the result of evaluating an expression as an integer or tuple. To bind these
    names,~\VC, like other functional logic languages, relies
    on~\it{unification} of its logical variables and \it{search} at runtime to
    meet a set of program constraints~\citep{antoy2010functional,
    hanus2013functional}. Combining unifying logical variables with search
    requires backtracking, which can lead to exponential runtime
    cost~\citep{hanus2013functional, wadler1985replace, clark1982introduction}. 

    % Todo: An example would be nice. Backtracking/evaluating backwards, then
    % done. 

    Pattern matching, by contrast, can be compiled to a~\it{decision tree}, a
    data structure that enforces linear runtime performance by guaranteeing no
    part of the scrutinee will be examined more than once~\citep{maranget}. A
    decision tree does not backtrack: once a program makes a decision based on
    the form of a value, it does not re-test it later with new information. 
    

    Figures: 

    \begin{figure}[ht] 
        \begin{minipage}[h]{0.54\linewidth}
          \verselst
          \begin{lstlisting}[numbers=none, basicstyle=\tiny, xleftmargin=.2em,
                            showstringspaces=false,
                            frame=single]
$\exists$ exclaimTall. exclaimTall = $\lambda$ sh. 
  one { 
      $\exists$ s. sh = $\langle$Square s$\rangle$; 
      s > 100.0; "Wow! That's a sizeable square!"
    | $\exists$ w h. sh = $\langle$Triangle, w, h$\rangle$; 
      h > 100.0; "Goodness! Towering triangle!"
    | $\exists$ b1 b2 h. sh = $\langle$Trapezoid, b1, b2, h$\rangle$;
      h > 100.0; "Zounds! Tremendous trapezoid!"
    | "Your shape is small." }
            \end{lstlisting}
                \Description{exclaimTall}
        \subcaption{\tt{exclaimTall} in~\VC }
            \label{fig:verseexclaimtall} 
        %   \vspace{4ex}
        \end{minipage}%%
        \begin{minipage}[h]{0.5\linewidth}
          \verselst
          \begin{lstlisting}[numbers=none, basicstyle=\tiny, xleftmargin=2em,
                        frame=single]
$\exists$ tripleLookup. tripleLookup = $\lambda$ rho x. 
    one { $\exists$ w. lookup rho x = $\langle$Just w$\rangle$; 
          $\exists$ y. lookup rho w = $\langle$Just y$\rangle$; 
          $\exists$ z. lookup rho y = $\langle$Just z$\rangle$;
          z
        | handleFailure x }
          \end{lstlisting}
                \Description{exclaimTall}
            \subcaption{\tt{tripleLookup} in~\VC }
            \label{fig:versetriplelookup} 
        \vspace{4ex}
        \end{minipage} 
        \begin{minipage}[h]{\linewidth}
          \verselst
          \begin{lstlisting}[numbers=none, basicstyle=\tiny, xleftmargin=9em, 
                            frame=single, showstringspaces=false]
$\exists$ game_of_token. game_of_token = $\lambda$ token. 
  $\exists$ f. one {
         token = one { $\langle$BattlePass, f$\rangle$ | $\langle$ChugJug, f$\rangle$ | $\langle$TomatoTown, f$\rangle$}; 
           $\langle$"Fortnite", f$\rangle$
       | token = one { $\langle$HuntersMark, f$\rangle$ | $\langle$SawCleaver, f$\rangle$}; 
           $\langle$"Bloodborne", 2 * f$\rangle$
       | token = one { $\langle$MoghLordOfBlood, f$\rangle$ | $\langle$PreatorRykard, f$\rangle$}; 
           $\langle$"Elden Ring", 3 * f$\rangle$
       |   $\langle$"Irrelevant", 0$\rangle$ }
          \end{lstlisting}
                \Description{game_of_token in~\VC}
        \subcaption{\tt{game\_of\_token} in~\VC }
            \label{fig:versegot} 
        \vspace{4ex}
        \end{minipage}%% 
    \caption{Code for the~\ref{pmandextensions} functions with equations looks
    similar, and it doesn't need extensions.}
    \label{fig:verseextfuncs}
      \end{figure}
        




\end{document}
\endinput
%%
